\documentclass[12pt,a4paper]{article}
        \makeatletter
        \def\input@path{{../}}
        \makeatother
        \usepackage{almeria-fichas}

        \FichaMeta{Sesion 11}{Variables urbanas}{Bloque 3 · Datos, funciones y graficas}{Identificar la variable independiente y la dependiente al analizar datos reales de la ciudad.}{Hechos}{Recordar, Comprender, Aplicar}

        \begin{document}

        \FichaCabecera

        \begin{materialesbox}
        \begin{itemize}
    \item Ficha impresa
    \item Lapiz
    \item Regla
    \item Calculadora si se necesita
\end{itemize}
        \end{materialesbox}

        \begin{pasosbox}
        \begin{enumerate}
    \item Lee la situacion y subraya que cambia primero.
    \item Marca con una x la variable de entrada y con una y la de salida.
    \item Haz una tabla si te ayuda a ordenar.
    \item Justifica una respuesta con una frase breve.
\end{enumerate}
        \end{pasosbox}

        \begin{recordatoriobox}
        \begin{itemize}
    \item Variable independiente: la que elegimos o controlamos.
    \item Variable dependiente: la que cambia segun la otra.
    \item En una tabla suele colocarse primero la variable independiente.
\end{itemize}
        \end{recordatoriobox}

        \begin{apoyobox}
        \begin{itemize}
    \item Pregunta guia: ¿de que depende este valor?
    \item Puedes usar las palabras entrada y salida.
    \item Si cambias el tiempo y cambia la temperatura, el tiempo es x.
\end{itemize}
        \end{apoyobox}

        \BloomSection{recordar}

\BloomExercise{recordar}{1}{Nombres de variables}{Define variable independiente y variable dependiente con una frase sencilla.}{5}

\BloomExercise{recordar}{2}{Entrada y salida}{Completa: en una tabla, la variable de entrada suele ser la \underline{\hspace{3cm}}.}{4}

\BloomSection{comprender}

\BloomExercise{comprender}{1}{Explica la dependencia}{En la situacion \emph{horas de aparcamiento y precio pagado}, explica cual es la variable independiente y cual la dependiente.}{6}

\BloomExercise{comprender}{2}{Otro ejemplo}{En \emph{mes del año y temperatura media}, explica por que una variable depende de la otra.}{6}

\BloomSection{aplicar}

\BloomExercise{aplicar}{1}{Clasifica parejas}{Indica las variables independiente y dependiente en estos casos: distancia-tiempo, numero de entradas-precio total, dia del mes-visitantes.}{7}

\BloomExercise{aplicar}{2}{Tabla simple}{Construye una tabla con una variable independiente y una dependiente a partir de una situacion real de Almeria.}{7}

        \begin{evidenciabox}
        \begin{itemize}
    \item Distinguir variable independiente y dependiente.
    \item Organizar datos en una tabla sencilla.
    \item Relacionar variables con contextos urbanos reales.
\end{itemize}
        \end{evidenciabox}

        \end{document}
